% Generated draft: requires mathematical and source review.
\documentclass[11pt]{article}
\usepackage[margin=25mm]{geometry}
\usepackage{fontspec}
\setmainfont{Latin Modern Roman}
\usepackage{amsmath,amssymb,mathtools,unicode-math}
\setmathfont{Latin Modern Math}
\usepackage{enumitem,needspace,xurl,hyperref}
\setlength{\parindent}{0pt}
\setlength{\parskip}{5pt}
\setlength{\emergencystretch}{3em}
\newcommand{\N}{\mathbb{N}}
\newcommand{\Z}{\mathbb{Z}}
\newcommand{\Q}{\mathbb{Q}}
\newcommand{\R}{\mathbb{R}}
\newcommand{\C}{\mathbb{C}}
\newcommand{\F}{\mathbb{F}}
\newcommand{\A}{\mathbb{A}}
\newcommand{\PP}{\mathbb{P}}
\newcommand{\E}{\mathbb{E}}
\newcommand{\eps}{\varepsilon}
\newcommand{\dd}{\,\mathrm{d}}
\begin{document}
\section*{UnsolvedMath ID: 30005943}
\textbf{Source problem code:} catalogue entry
\subsection{Definitions and mathematical statement}
Let \(X\) be a real target in the category intended by the supplied report. Real Gromov--Witten theory concerns curve-counting invariants compatible with the real structures on the target and domain. The source assumes that \(X\) is real-orientable and that its quantum cohomology is semisimple. The precise category of targets, real-orientation data, coefficient ring, and meaning of semisimplicity at the relevant quantum parameters are not specified in the supplied excerpt; these hypotheses must be recovered from the report before a complete formulation is possible.

Write \(t,q\) for the parameters used in the source and \(\psi\) for a cotangent-line class at a marked point. A descendant insertion \(\psi^{k_i}\gamma_i\) consists of a nonnegative power of this class and a cohomology insertion \(\gamma_i\). The excerpt places the proposed identity in the range
\[
1-g-n<0,\qquad \tau(t)=t,
\]
where \(g\) is the genus and \(n\) the number of displayed descendant insertions. The action denoted by \(\tau\), the permitted insertions, and the parameter domains are not defined in the supplied record.

\textbf{Conjecture 1, as recoverable from the supplied record.} For every real-orientable target \(X\) whose quantum cohomology is semisimple, there exists a suitable series \(U_{t,q}(\psi)\) such that the reconstruction identity called Theorem 1 in the cited report holds for its real Gromov--Witten theory, throughout the stated range \(1-g-n<0\) and at every permitted point \(t\) fixed by \(\tau\). The intended identity reconstructs the relevant real descendant invariants from integrals of \(\psi\)-classes over moduli spaces of curves and the genus-zero complex and real Gromov--Witten theory of \(X\).

The excerpt describes a graph sum with weights \(1/|\operatorname{Aut}(\Gamma)|\). Its vertex forms have \(n_v+k\) arguments, indexed by incident half-edges, marked leaves, real leaves, and \(k\) additional insertions. The insertion rules involve \(R_{t,q}\), \(S_{t,q}\), \(E_{t,q}\), \(T_{t,q}\), involution operators \(\widehat\tau,\widehat\tau'\), and coefficients of \(z^{-k_i-1}\). Transposes are taken with respect to the Poincar\'e pairing; the stated real-leaf insertion is \(U_{t,q}(-\psi)^T\). However, the displayed theorem and several insertion formulas are corrupted in the supplied extraction. The graph class, normalization, operator definitions, and type and admissibility conditions of \(U_{t,q}\) cannot be recovered reliably. Thus this is a source-dependent statement of the conjectured extension, not a complete reconstruction formula.

The supplied dated curation says that the report proves the identity for \((\mathbb P^{2N-1},\theta_N)\) with its stated involution, while the extension to every target satisfying the conjectural hypotheses remains the intended question. The involution \(\theta_N\) is not defined in the supplied excerpt. This reported special case and the supplied status label are not independently verified here.
\subsection{Short English statement}
The conjecture asks whether a suitable series \(U_{t,q}(\psi)\) extends the report's reconstruction formula from its stated real projective targets to every real-orientable target with semisimple quantum cohomology. The supplied range is \(1-g-n<0\), with \(\tau(t)=t\). Corruption of the extracted formula and missing definitions prevent a fully explicit identity; human review against the cited report is required.
\subsection{Sources}
\begin{itemize}
\item Local UnsolvedMath registry, permanent ID \(30005943\), source problem code \texttt{catalogue entry}, title: \emph{Reconstruction of Real Gromov--Witten Invariants for Semisimple Targets}. Supplied collection: Oberwolfach Reports Open Problems; proposed year: 2024. No individual proposer is supplied.
\item Supplied report citation: \emph{Real and Logarithmic Enumerative Geometry} (2024), Oberwolfach Reports, pp.~1285--1338, DOI 10.4171/owr/2024/23. Source URL and supplied literature reference labelled ``Oberwolfach Report 23/2024'': \url{https://doi.org/10.4171/owr/2024/23}.
\item Literature assessment dated exactly \texttt{2026-08-22}: the supplied curation labels the problem ``partially solved'' and states that the report proves the reconstruction identity for the real projective targets with the stated involution, whereas Conjecture 1 concerns all real-orientable targets with semisimple quantum cohomology. This is supplied dated curation derived from the OWR record, not an independently verified literature assessment.
\item The registry labels the statement \texttt{corrected\_verified} and reports that it was checked against the official report. That verification claim is supplied provenance; the report has not been consulted for this draft, and the supplied original-statement extraction remains incomplete and corrupted.
\end{itemize}
\end{document}
